% Generated draft: requires mathematical and source review.
\documentclass[11pt]{article}
\usepackage[margin=25mm]{geometry}
\usepackage{fontspec}
\setmainfont{Latin Modern Roman}
\usepackage{amsmath,amssymb,mathtools,unicode-math}
\setmathfont{Latin Modern Math}
\usepackage{enumitem,needspace,xurl,hyperref}
\setlength{\parindent}{0pt}
\setlength{\parskip}{5pt}
\setlength{\emergencystretch}{3em}
\newcommand{\N}{\mathbb{N}}
\newcommand{\Z}{\mathbb{Z}}
\newcommand{\Q}{\mathbb{Q}}
\newcommand{\R}{\mathbb{R}}
\newcommand{\C}{\mathbb{C}}
\newcommand{\F}{\mathbb{F}}
\newcommand{\A}{\mathbb{A}}
\newcommand{\PP}{\mathbb{P}}
\newcommand{\E}{\mathbb{E}}
\newcommand{\eps}{\varepsilon}
\newcommand{\dd}{\,\mathrm{d}}
\begin{document}
\section*{UnsolvedMath ID: 30003246}
\textbf{Source problem code:} catalogue entry
\subsection{Definitions and mathematical statement}
Let \(j\) denote the modular \(j\)-function referred to in the source, with \(j'\) and \(j''\) denoting its first and second derivatives in the intended analytic variable. The source uses \(j\operatorname{cl}\) for a closure associated with \(j\)-derivations, but does not supply its definition, its ambient structure, or the definitions of the transformation groups \(G\) and \(G_F\). These omissions remain unresolved here.

Retain the supplied hypotheses: \(C\) is \(j\operatorname{cl}\)-closed, \(\lambda\) is a \(j\operatorname{cl}\)-independent tuple, and \(F=\Q(\lambda)\). Let \(z=(z_1,\ldots,z_n)\) be a tuple whose coordinates lie in distinct \(G\)-orbits, and let \(g\in G_F\setminus G\) be nonscalar. The domain of \(z\), the action defining \(gz\), and the conditions ensuring that all displayed function values exist must be supplied by the intended ambient setting; they are not specified in the record. It is also unspecified whether independence of \(\lambda\) is over \(C\) or over another base.

For a tuple \(w=(w_1,\ldots,w_n)\) in the domain of these functions, write
\[
J(w)=\bigl(j(w_i),j'(w_i),j''(w_i)\bigr)_{i=1}^{n}.
\]
In the intended common ambient field, let \(K\) be the smallest subfield containing \(F\cup C\). Interpret the source notation \(\operatorname{td}(J(z),J(gz)/F,C)\) as
\[
\operatorname{trdeg}_{K} K\bigl(J(z),J(gz)\bigr),
\]
where transcendence degree is the cardinality of a transcendence basis of the indicated field extension.

The supplied problem asks whether the bound attributed to the source's main theorem,
\[
\operatorname{trdeg}_{K}K\bigl(J(z),J(gz)\bigr)\geq 3n,
\]
is sharp under the intended hypotheses, or whether a stronger lower bound can be established without proving a modular Schanuel conjecture. It also asks for an explicit description of the \(j\)-derivation pregeometry \(j\operatorname{cl}\). The record does not specify the parameter quantifiers intended for sharpness, the proposed form of an improvement, the particular modular Schanuel conjecture meant, or what constitutes the requested explicit description. Thus this is a statement of the supplied research targets with unresolved hypotheses, not a complete independently verified theorem or conjecture.
\subsection{Short English statement}
The problem asks whether the stated \(3n\) transcendence-degree bound for modular \(j\)-values and their first two derivatives is optimal, whether it can be strengthened without proving a modular Schanuel conjecture, and how to describe \(j\operatorname{cl}\) explicitly. The supplied record leaves essential group, domain, closure, and independence conventions unspecified.
\subsection{Sources}
\begin{itemize}
\item Supplied local UnsolvedMath registry record: permanent ID \(30003246\), exact source problem code \texttt{catalogue entry}, title \emph{Sharp Transcendence Bounds for the Modular J-Function}, collection \emph{Oberwolfach Reports Open Problems}. The proposed year is recorded as 2016; no individual proposer is supplied.
\item Source citation supplied by the registry: \emph{Definability and Decidability Problems in Number Theory} (2017), Oberwolfach Reports, pp.~2793--2866, DOI \texttt{10.4171/owr/2016/49}. \url{https://doi.org/10.4171/owr/2016/49}. The background describes Oberwolfach Reports as workshop reports documenting open-problem discussions. No theorem number or problem-specific page is supplied.
\item Supplied literature reference: Eterovi\'c, \emph{Generic solutions}. \url{https://doi.org/10.1007/s00208-024-03082-6}.
\item Supplied literature reference: \emph{Extending j-derivations}. \url{https://ueaeprints.uea.ac.uk/80428/1/Extending_j_derivations.pdf}.
\item Supplied literature reference: \emph{Oberwolfach Reports 2017}. \url{https://ems.press/content/serial-article-files/46652}.
\item Literature assessment dated exactly \texttt{2026-08-21}: the supplied curation labels the problem partially solved and reports conditional and structural advances, including a pregeometry result, while stating that neither the sharp inequality nor the requested full explicit description is established. This is supplied dated curation derived from the OWR record, not an independently verified literature assessment; the status and cited claims require independent verification.
\item The registry labels the statement \texttt{corrected\_verified} and reports that it was checked against the official Oberwolfach report, preserving that report's hypotheses, notation, and target. This is a supplied verification claim, not verification performed for this draft.
\end{itemize}
\end{document}
